% =====================================================================
% Green IT 2026 (MICS2-63): Peer Review 1
%
% One file per paper reviewed. You submit TWO: this one and
% peer-review-2.tex. They are identical templates; fill each in against
% a different paper.
%
% Length: ~2 double-column pages each.
%
% Build:  latexmk        (from the folder above this one, not from src/)
% =====================================================================
\documentclass[conference,a4paper]{IEEEtran}

\usepackage[T1]{fontenc}
\usepackage{graphicx}
\usepackage{booktabs}
\usepackage{array}
\usepackage{url}
\usepackage{amsmath}

% ---------------------------------------------------------------------
% Guidance notes. Each one tells you what its section owes.
% Delete a note as you write its section, and set \guidancetrue below to
% \guidancefalse before you submit: the PDF you hand in carries no guidance.
\newif\ifguidance
\guidancetrue
\newcommand{\guidance}[1]{\ifguidance{\itshape\small #1\par\medskip}\fi}

\begin{document}

\title{Peer Review 1: \textit{<title of the paper you are reviewing>}}

\author{\IEEEauthorblockN{Anonymous reviewer}
\IEEEauthorblockA{Green IT 2026 (MICS2-63)\\
Master in Information and Computer Sciences,
University of Luxembourg\\[2pt]
% Build stamp: the date this file was last compiled.
\scriptsize Compiled \today.}}

\maketitle


\guidance{Every italic note in this document is guidance.
Delete each one as you write its section, then set
\texttt{\textbackslash guidancefalse} in the preamble: the PDF you submit
carries no guidance.

The sections below are the parts of the paper, in the paper's own order. Write
about each one where it belongs: a defect goes in the section it affects.

Reviews are double-blind. Your name stays out of this document. Your
submission on Moodle is what tells the instructor the review is yours. The
instructor numbers the reviews when he forwards them, from the author's side,
so this file being your first review says nothing about the number it arrives
under.

The paper you receive comes without its AI Usage Journal, which the author
wrote and the instructor marks separately, and with an archive holding the
script you are asked to run in Section~IV.

This review carries no score of any kind. Each author receives two reviews of
their paper. Marks reach students in mid-February.}

% ---------------------------------------------------------------------
\section{Introduction}
\guidance{What question does the paper set out to answer, and does the paper
answer it? Take the claims the introduction makes and say whether the rest of
the paper supports them.}

% ---------------------------------------------------------------------
\section{Related Work}
\guidance{Check the references. Do they exist, and do they say what they are
cited for? Then judge whether the section situates the author's own
contribution, or only lists other people's.}

% ---------------------------------------------------------------------
\section{Life Cycle Assessment}
\guidance{Judge each subsection of III on whether it does its job. Then
recompute the baseline from what the paper gives you, and put your figure in
Table~\ref{tab:recomputation}. Section~III.B declares what makes the paper's
numbers reproducible: for a web page, whether it arrives compressed and on what
date that was read, so recompute the bytes on that basis; for an HPC job, the
commit, the fixed input and the exclusive node it ran on. If you cannot recompute it, say
precisely what is missing.}

% ---------------------------------------------------------------------
\section{Optimisation and Payback}
\guidance{Run the script from the archive distributed with the paper and report
what happened. Keep a scientific finding apart from an environment problem.

Recompute the remaining figures of Table~\ref{tab:recomputation}. A before and
an after are comparable when they are produced the same way: for a web page,
the same server started with the same compression setting; for an HPC job, the
same partition and node type with the same fixed input. Check that Section~IV
says which. Then judge whether the optimisation and the payback follow from
what Section~III found.}

\begin{table}[htbp]
\centering
\caption{Recomputation of the author's key figures.}
\label{tab:recomputation}
\begin{tabular}{|>{\raggedright\arraybackslash}p{2.3cm}|r|r|>{\raggedright\arraybackslash}p{2.0cm}|}
\hline
\textbf{Quantity} & \textbf{Author} & \textbf{Reviewer} & \textbf{Comment} \\
\hline
Baseline, live page or job &  &  &  \\
\hline
Before optimisation, frozen copy or original job &  &  &  \\
\hline
After optimisation       &  &  &  \\
\hline
Development cost         &  &  &  \\
\hline
Payback (page views or runs) &  &  &  \\
\hline
\end{tabular}
\end{table}

% ---------------------------------------------------------------------
\section{Threats to Validity}
\guidance{Judge the threats the author lists, then name one they did not. If
you looked and found none, say where you looked.}

% ---------------------------------------------------------------------
\section{Recommendation}
\guidance{One word: ACCEPT or REJECT. Those are the two options. Taking one
of them is what this section is for.

Then two or three lines saying why this verdict follows from what you wrote
above. Which side you take is yours.}

% ---------------------------------------------------------------------
\section{Actionable Recommendations}
\guidance{For each issue you raised above, what the author should do. Concrete
enough to carry out before Fri 18 Dec.}

\end{document}
